\section{问题重述}
某微网由小区负荷、光伏发电、储能装置及外部电网构成。外部电网负责补足本地供能缺口，储能可在低价或光伏富余时充电，在高价或供电不足时放电。题目给出固定日内电价与典型日曲线、全年负荷和光伏实测数据、每日四次光伏预报以及全年波动电价，要求制定经济且可执行的购电策略。

问题一要求在负荷、光伏和电价均已知时安排144个10分钟时段的购电与储能动作；问题二要求在每日0:00仅凭历史信息制定全天计划，并以5倍电价补充紧急购电；问题三允许在6:00、12:00和18:00依据新增预报调整尚未执行的购电量；问题四用波动电价分别重算问题二和问题三。各问共享储能物理约束，但信息集和结算方式逐层扩展。

本文总体技术路线见图~\ref{fig:roadmap}。核心原则是先统一数据口径和物理可行域，再逐层增加预测误差、滚动更新与价格波动，最后以独立回代和简单基线检验模型收益。
\begin{figure}[H]\centering
\includegraphics[height=0.78\textheight]{../figures/fig_roadmap.pdf}
\caption{微网购电与储能协同调度总体技术路线}\label{fig:roadmap}
\end{figure}
